
\paragraph{Realizar valoración}
En la figura \ref{FIG:VAR}, se observa el comportamiento del sistema cuando el usuario realiza la valoración del videojuego. 
Una vez que el usuario se autentica, este podrá ver la información del videojuego y, a continuación, realizar la valoración. 
Para llevar a cabo la valoración, debe rellenar el formulario. Una vez que el usuario envía el formulario, el backend valida 
la información y responde 'OK' que indica que la operación ha sido exitosa.

\begin{figure}[Flujo realizar valoración]{FIG:VAR}{En este diagrama se puede ver el proceso completo de la valoración de un videojuego}
    \image{10cm}{}{valorar}
\end{figure}

\paragraph{Crear lista}
En la figura \ref{FIG:VAR}, se observa el comportamiento del sistema cuando el usuario crea una lista personalizada. Tras autenticarse, 
el usuario accede a la sección de 'Mis listas' y selecciona la opción de 'Crear lista'. A continuación, rellena el formulario con la información adecuada
que es enviada al backend, quien almacena la nueva lista y devuelve 'OK' que indica que la operación ha sido exitosa.

\begin{figure}[Flujo crear lista personalizada]{FIG:VAR2}{En este diagrama se puede ver el proceso completo de la creación de la lista personalizada }
    \image{10cm}{}{crearlista}
\end{figure}

\paragraph{Editar lista}
Tras completar el paso anterior, el usuario puede editar su lista. El usuario selecciona la lista que desea
actualizar. A continuación, rellena el formulario con los datos nuevos. Esta información es enviada al backend, que actualiza el contenido 
y responde 'OK'. Este proceso se ilustra en la figura \ref{FIG:VAR3}

\begin{figure}[Flujo editar lista personalizada]{FIG:VAR3}{En este diagrama se puede ver el proceso completo de la edición de la lista personalizada }
    \image{10cm}{}{editarlista}
\end{figure}
